
Section~\ref{sec:practice} states what the pilot is and what it is not. This
appendix gives the flow, the design and the numbers.

\subsection*{Flow}

\begin{table}[t]
\centering
\small
\resizebox{\ifdim\width>\linewidth\linewidth\else\width\fi}{!}{%%
\begin{tabular}{lr}
\toprule
step & n \\
\midrule
records enumerated in the frame & 54910.000 \\
records sampled & 600.000 \\
full text retrieved & 369.000 \\
full text not retrieved & 231.000 \\
screened in by the mechanical rule & 54.000 \\
screened out: no quantitative metric & 192.000 \\
screened out: no incremental comparison & 123.000 \\
adjudicated, screened-in stratum (census) & 54.000 \\
adjudicated, screened-out stratum (sample) & 120.000 \\
confirmed eligible, screened-in stratum & 26.000 \\
confirmed eligible, screened-out stratum (unweighted) & 11.000 \\
estimated eligible among all retrieved full texts & 54.900 \\
\bottomrule
\end{tabular}%
}
\caption{The pilot's flow, in PRISMA-style steps.}
\label{tab:prisma}
\end{table}


Table~\ref{tab:prisma} is the flow, laid out in the style
\citet{page2021prisma} prescribe for a systematic review --- which this is
not, and the appendix says so throughout; the layout is used because it makes
what was excluded, and why, countable. Records were enumerated from
a public index across \nAuditStrata\ strata --- venues, two citation seeds and
a topic frame --- and the enumeration is cached per stratum in the archive so
it can be re-run without re-querying.

The allocation was declared before the enumeration finished: proportional to
stratum size with a floor of \minPerStratum\ per stratum, up to a budget of
\nAuditTarget\ records, with every estimate carrying the design weight
$N_h / n_h$. The enumeration returned \nAuditFrame\ deduplicated records ---
short of that budget, because the provider's daily request quota was exhausted
first --- so the allocation bound at $n_h = N_h$ in every stratum and the
design weight is \auditMaxWeight\ throughout. \textbf{The machinery of
weighting is therefore inert here.} It is retained because it is what makes
the design re-runnable at scale, and because a reader is entitled to see that
the weights are one rather than to assume it; but the pilot's reach is that of
a census of \nAuditFrame\ enumerated records, not that of a probability sample
of a literature, and Section~\ref{sec:practice} says so.

\subsection*{The two-stratum adjudication, and why it is the point}

\begin{table}[t]
\centering
\small
\resizebox{\ifdim\width>\linewidth\linewidth\else\width\fi}{!}{%%
\begin{tabular}{lrrrl}
\toprule
quantity & estimate & lo & hi & note \\
\midrule
screen sensitivity & 0.474 & 0.347 & 0.602 & weighted; the screened-out stratum is a 38.1\% sample \\
screen specificity & 0.911 & 0.874 & 0.938 & weighted \\
screen precision (PPV) & 0.481 & 0.354 & 0.611 & census over the screened-in stratum \\
eligibility prevalence among retrieved full texts & 0.149 & -- & -- & two-stratum estimate \\
\bottomrule
\end{tabular}%
}
\caption{The mechanical screen's measured accuracy against the adjudicated standard.}
\label{tab:screenacc}
\end{table}


The design's one non-obvious feature is that the papers the mechanical screen
\emph{rejected} were sampled and adjudicated too. Without that, a screen's
misses are invisible and its proportions can only be described as bounds ---
which an earlier version of this work did, incorrectly, because a screen with
false positives does not yield a bound in either direction.

With it, Table~\ref{tab:screenacc} reports what the screen actually does:
sensitivity \auditSens\ [\auditSensLo, \auditSensHi], specificity
\auditSpec, precision \auditPrec. Agreement between the mechanical codes and
the adjudicated ones has median $\kappa = \auditKappa$, which is poor. The
prevalence estimates in Table~\ref{tab:pilot} are nonetheless unbiased for the
adjudicated standard, because the two strata are combined with their sampling
weights rather than the screened-in stratum being used alone. That is why no
sensitivity-and-specificity correction of the \citet{rogan1978estimating} kind
is applied: such a correction repairs a prevalence estimated \emph{from the
screen}, and the estimate here is made from the adjudications, with the screen
entering only as the stratifier.

\subsection*{Two mechanical screens, one standard}

\begin{table}[t]
\centering
\small
\resizebox{\ifdim\width>\linewidth\linewidth\else\width\fi}{!}{%%
\begin{tabular}{lrrrr}
\toprule
rule & sensitivity & specificity & precision & n flagged \\
\midrule
amended (primary) & 0.474 & 0.911 & 0.481 & 54 \\
registered & 0.569 & 0.769 & 0.301 & 73 \\
\bottomrule
\end{tabular}%
}
\caption{Two mechanical screens against the same adjudicated standard. They fail differently and neither is good enough to support a claim about a field's practice, which is why the pilot does not make one.}
\label{tab:twoscreens}
\end{table}


Two screening rules ran on every paper: the one registered in
\texttt{AUDIT-PROTOCOL.md} section~4, and the amended one of its
amendment~1, written against the same prose. Table~\ref{tab:twoscreens}
reports both against the adjudicated standard. They fail differently --- the
registered rule flags more papers and is right about fewer of them --- and
neither is good enough to support a claim about a field's practice. That is
the pilot's central finding about itself.

\subsection*{Applicability-specific denominators}

\begin{table}[t]
\centering
\small
\resizebox{\ifdim\width>\linewidth\linewidth\else\width\fi}{!}{%%
\begin{tabular}{lr}
\toprule
quantity & n \\
\midrule
eligible papers & 54.875 \\
of which the incremental feature is a lookup into an external register or reference source & 3.000 \\
of those, reporting the increment across a range of register populations & 0.000 \\
\bottomrule
\end{tabular}%
}
\caption{Applicability-specific denominators.}
\label{tab:applicability}
\end{table}


Not every code applies to every eligible paper. The register-population code
is meaningful only where the incremental feature is a lookup into an external
register or reference source whose per-case coverage can be incomplete, and
Table~\ref{tab:applicability} reports that denominator separately. Of
\nAuditEligible\ estimated eligible papers, \nAuditApplicable\ have such a
feature. A proportion of zero over a denominator of \nAuditApplicable\ is a
much weaker statement than a proportion of zero over \nAuditEligible, and the
earlier version reported the latter.

\subsection*{What is missing, and what it could do}

\begin{table}[t]
\centering
\small
\begin{tabular}{lrr}
\toprule
characteristic & retrieved & not\_retrieved \\
\midrule
n & 369.000 & 231.000 \\
median year & 2023.000 & 2024.000 \\
share from the four round-eighteen venues & 0.650 & 0.654 \\
\bottomrule
\end{tabular}

\caption{Retrieved against unretrieved records, and the bounds on the eligibility rate under the extreme assumptions about the unretrieved.}
\label{tab:missing}
\begin{tabular}{lr}
\toprule
assumption & p\_eligible \\
\midrule
none of the unretrieved papers is eligible & 0.091 \\
the unretrieved papers are eligible at the retrieved rate & 0.149 \\
all of the unretrieved papers are eligible & 0.476 \\
\bottomrule
\end{tabular}
\end{table}


Full text was unavailable for \nAuditNoFullText\ of the sampled records.
Table~\ref{tab:missing} compares the retrieved and the unretrieved on every
characteristic the metadata carries, and gives the eligibility rate under the
extreme assumptions that none and that all of the unretrieved are eligible:
\auditBoundLo\ to \auditBoundHi. The retrieved and unretrieved sets do not
differ materially on year or venue mix, which is reassuring and is not
evidence of exchangeability.

\subsection*{What this design can and cannot resolve}

\begin{table}[t]
\centering
\small
\resizebox{\ifdim\width>\linewidth\linewidth\else\width\fi}{!}{%%
\begin{tabular}{rll}
\toprule
half width & n required & achieved \\
\midrule
0.050 & 385 & 29.6 \\
0.070 & 196 & 29.6 \\
0.100 & 97 & 29.6 \\
0.150 & 43 & 29.6 \\
0.200 & 25 & 29.6 \\
\bottomrule
\end{tabular}%
}
\caption{What the pilot's design can and cannot resolve: the adjudicated sample each half-width would need, against the effective sample the pilot achieved.}
\label{tab:power}
\end{table}


An effective sample of \nAuditEffective\ resolves a proportion near one half
to about $\pm\auditHalfWidthPct$ percentage points. A design that resolved it
to $\pm 7$ would need \nAuditRequired\ adjudicated eligible papers.
Table~\ref{tab:power} gives the requirement at five half-widths. This is the
single most important limitation of the pilot and it is why no claim in the
paper depends on it.

\subsection*{Adjudication materials}

Every adjudication was made from an evidence dossier extracted from the paper's
full text: the sentences the screen fired on, the sentences carrying a
performance metric beside a number, and the neighbourhood of every mention of
each coded concept. The dossiers are committed to
\texttt{data/audit2/dossiers/} and the judgements to
\texttt{data/audit2/adjudication\_in.csv} and
\texttt{adjudication\_out.csv}, with a free-text reason for every one. There
is one adjudicator, machine-assisted; that is transparency, not independence,
and Section~\ref{sec:practice} says so.

\subsection*{The prevalence estimates, and their denominators}

\begin{table}[t]
\centering
\small
\resizebox{\ifdim\width>\linewidth\linewidth\else\width\fi}{!}{%%
\begin{tabular}{lrrrrrrr}
\toprule
code & estimate & boot lo & boot hi & design lo & design hi & Wilson lo & Wilson hi \\
\midrule
B\_stated & 0.850 & 0.705 & 0.973 & 0.755 & 0.945 & 0.682 & 0.937 \\
B\_justified & 0.175 & 0.067 & 0.315 & 0.102 & 0.249 & 0.078 & 0.347 \\
M\_justified & 0.055 & 0.000 & 0.129 & 0.042 & 0.067 & 0.013 & 0.198 \\
Theta\_stated & 0.271 & 0.121 & 0.442 & 0.162 & 0.380 & 0.144 & 0.450 \\
Range\_reported & 0.440 & 0.264 & 0.619 & 0.317 & 0.562 & 0.278 & 0.615 \\
\bottomrule
\end{tabular}%
}
\caption{The pilot's proportions over the estimated eligible population, with the stratified bootstrap interval this paper reports, the design-based linearised interval beside it, and the Wilson interval an earlier version printed. None of the five Wilson intervals is contained in its design-based counterpart.}
\label{tab:pilot}
\end{table}


Over an estimated \nAuditEligible\ eligible papers with an effective sample
size of \nAuditEffective: \auditBStatedPct\ state the composition of their
baseline; \auditBJustifiedPct\ argue for what the baseline excludes;
\auditMJustifiedPct\ argue for the metric; \auditThetaStatedPct\ state or
integrate over an operating point; \auditRangeBaselinePct\ report the
increment at more than one baseline, \auditRangeMetricPct\ under more than one
metric, \auditRangeThresholdPct\ across a range of operating points and
\auditRangePopulationPct\ across a range of register populations.

Denominators are applicability-specific. Only \nAuditApplicable\ of the
eligible papers have an incremental feature that is a lookup into an external
register at all, so the population axis has a denominator of
\nAuditApplicable\ and not of \nAuditEligible. Using the eligible count for
every axis understates every rate whose axis does not apply to every paper,
and is not done here.

\paragraph{Variance}
Each proportion here is a \emph{ratio} of two weighted totals: the numerator
counts eligible papers carrying a code, the denominator estimates the eligible
population, and both are random. The screened-in stratum is a census, so it
carries a design weight of \auditWeightIn\ and contributes no sampling
variance; the screened-out stratum was adjudicated at a sampling fraction
whose reciprocal is \auditWeightOut.

A \textbf{Wilson} interval computed from a fractional ``effective sample
size'' --- which is what this pilot previously reported --- is a binomial
interval for a binomial count, and neither the numerator nor the denominator
here is one. Table~\ref{tab:pilot} reports instead a
\textbf{stratified nonparametric bootstrap} --- resampling adjudicated papers
within stratum and recomputing the ratio, \nPilotBoot\ times --- as the
primary interval, because with \nAuditOutSampled\ adjudicated papers in the
second stratum a linearisation is not to be trusted, and the
\textbf{design-based} linearised interval beside it. The design-based interval
has a median width of \auditDesignWidthPct\ percentage points and the
bootstrap \auditBootWidthPct, against the Wilson interval's
\auditWilsonWidthPct; and \emph{none} of the five Wilson intervals is
contained in its design-based counterpart, the endpoints differing by up to
\auditMaxGapPct\ percentage points. That is correction~C28 in
\ref{app:corrections}.

\subsection*{Four limits, stated rather than buried}

The frame is what one provider's index returned before its daily quota stopped
the enumeration, so it is a convenience frame and its coverage of the field is
unmeasured; nothing here corrects for that, and nothing can without a second
index.

An effective sample of \nAuditEffective\ resolves a proportion to about
$\pm\auditHalfWidthPct$ percentage points; a design that resolved it to
$\pm 7$ points would need \nAuditRequired\ adjudicated papers. That is the
sample-size calculation, and this pilot does not meet it.

Full text was unavailable for \nAuditNoFullText\ of the sampled records, and
under the extreme assumptions that none or all of them are eligible the
eligibility rate moves between \auditBoundLo\ and \auditBoundHi.

There is one adjudicator, machine-assisted, not two independent human raters.
Every adjudication was made from an evidence dossier committed to the
repository, so the judgements can be re-examined against exactly what they
were made from; that is transparency, not independence, and no inter-rater
reliability statistic reported here is a substitute for a second human.
